\Needspace{14\baselineskip}
\section{Discussion}
\label{sec:discussao}

% PRISMA Item 17: Discussion of results in the context of other evidence
% PRISMA Item 18: Limitations of the review

\subsection{Main findings}
\label{sec:principais-achados}

The 28 mapped studies reveal a field concentrated around one dominant
architecture: experience banks that distil resolution trajectories into
reusable entries (46\% of primary studies,
\autoref{tab:arquiteturas}), with the LLM itself acting as the memory
controller (54\%) and evaluation centred on SWE-bench Verified (67\% of
studies with a benchmark, \autoref{tab:benchmarks}).
This convergence, however, does not imply maturity.
The field remains dominated by preprints (82\%), and the 28 studies use
at least eight distinct LLMs as backbones.
Controlled comparisons across architectures on the same benchmark with
the same model are absent from the corpus, so the observed convergence
reflects independent author choices rather than an empirical validation
of experience banks over alternatives such as knowledge graphs or Git
metaphors.

All 19 studies with controlled comparison report improvements over the
no-memory baseline (\autoref{tab:desempenho}), with a median gain of
6.8~pp on SWE-bench Verified and an amplitude of +2.2~pp to +39.0~pp.
The data indicate a diminishing-returns pattern:
on SWE-bench Verified, systems with baselines below 35\% obtain gains
above
10~pp~\citep{li2026traversalaspolicy,zhang2025darwingodel},
while baselines above 65\% obtain gains of 2 to
9~pp~\citep{deng2026memcoder,wu2025gcc,chen2025sweexp}.
The cost impact is bimodal: systems that replace context with memory
cut tokens by up to
39\%~\citep{li2026traversalaspolicy}.
Systems that add context via retrieval raise cost by +5 to
+21\%~\citep{chen2025sweexp,wang2026memgovern,xia2025livesweagent}.

The risk of degradation is not negligible.
Nine of the 28 studies (32\%) report scenarios where memory hurts
performance (\autoref{tab:degradacao}).
Retrieval noise is the mechanism cited in at least three of the nine
studies with degradation:
\citet{lindenbauer2025ctimrover} document a drop of 11~pp from unfiltered
item injection into context,
\citet{chen2025sweexp} document degradation with $k > 1$ retrieved
experiences, and
\citet{deshpande2025memtrack} identify that commercial memory backends
(Mem0, Zep) do not improve performance and keep tool-call redundancy at
roughly 21\%, with no reduction relative to the no-memory baseline.
The poisoning vulnerability documented by
\citet{srivastava2025memorygraft}, in which 9.1\% of malicious
experiences compromise 47.9\% of retrievals under dual-channel
retrieval (lexical + vector), indicates that memory stores without
provenance checks can constitute an attack surface.

\subsection{Implications for research}
\label{sec:implicacoes-pesquisa}

The central gap identified by this review is the absence of head-to-head
comparisons.
None of the 28 studies evaluates three or more persistence architectures
on the same benchmark, with the same model, and under the same compute
budget.
Each system reports gains against its own no-memory baseline, but
experimental conditions differ in model, framework, instance subset, and
number of attempts.
This fragmentation prevents any conclusion on which architecture is
superior, because the observed gains may reflect implementation
differences as much as differences in memory architecture.
A comparative study evaluating at least three representative
architectures under controlled conditions would fill this gap.

Cost underreporting likewise limits the practical utility of the
results.
Only seven studies provide a complete
breakdown~\citep{chen2025sweexp,guo2026eet,wang2026memgovern,wu2025gcc,xia2025livesweagent,zhang2025accelopt,zhou2025tomswe},
and 12 of the 28 (43\%) report no cost data at all.
Without standardised cost-per-task metrics (tokens consumed, monetary
cost, wall-clock time), practitioners cannot assess the trade-off of
adopting persistence.
Constructing a Pareto frontier between accuracy and cost, absent in
every analysed study, is a research direction with immediate applied
value.

Evidence on the relationship between architectural complexity and gain
is partial.
The data indicate that simple systems, such as
AccelOpt~\citep{zhang2025accelopt} with an experience queue and
EET~\citep{guo2026eet} with heuristic early termination, obtain
relevant gains in their respective domains without resorting to
multi-tier architectures.
However, this comparison is between distinct studies: benchmarks and
models differ across systems.
Only a study controlling the complexity variable internally could
confirm whether architectural simplicity is indeed sufficient or
whether scenarios exist where sophistication is justified.

Three studies propose memory as a directly executable artefact, an
emerging trend with still-limited evidence:
behaviour trees~\citep{li2026traversalaspolicy},
scripts created at run time~\citep{xia2025livesweagent}, and
whole agents~\citep{zhang2025darwingodel}.
The first two report substantial gains (+39.0~pp and +3.2~pp,
respectively), but the certainty of this finding is low given the
small number of studies and the absence of independent replication.

\subsection{Implications for practice}
\label{sec:implicacoes-pratica}

For teams that wish to incorporate knowledge persistence into coding
agents, the data in this review suggest a conservative starting point:
experience banks with similarity-based retrieval.
This is the dominant architecture (46\%), has reference open-source
implementations, and yields gains of 2 to 9~pp in every study with
baselines above 65\%.
More complex systems (knowledge graphs, hierarchical memory) do not
present, in the available data, clear evidence of systematically larger
gains, though this comparison is limited by benchmark and model
differences across studies.

The available data suggest that retrieval-mechanism calibration may
matter more than architectural sophistication.
\citet{chen2025sweexp} report that retrieving more than one experience
($k > 1$) degrades performance from 42.0\% to 39.0\%, and
\citet{shen2026structurally} report that category isolation outperforms
global retrieval by 2.3~pp.
Both results point in the same direction: selectivity in retrieval
influences the effective gain, and excess retrieved context can
introduce noise.

For already-deployed systems, experience-guided early termination
offers cost reduction with minimal risk.
\citet{guo2026eet} report 31.8\% savings in total cost with a maximum
loss of 0.2~pp, by identifying termination opportunities in 11.3\% of
instances.
This approach is agent-agnostic and requires no changes to the
underlying memory architecture.

Of the 24 primary studies, 17 share a complete repository and three
offer partial access, totalling 83\% with some degree of
reproducibility.

\subsection{Limitations of this review}
\label{sec:limitacoes}

The limitations of this review are organised into four categories of
validity threat: internal validity (rigour of the selection and
extraction processes), external validity (generalisation of the
findings beyond the corpus), construct validity (operationalisation of
the measured concepts), and conclusion validity (stability of
cross-study inferences).

Screening was conducted by a single researcher, which is an
acknowledged threat in systematic reviews and precludes the computation
of formal agreement indices such as Cohen's $\kappa$.
This limitation was mitigated by eligibility criteria calibrated in a
pilot study and refined through a stress test against five comparable
systematic reviews in
LLM/SE~\citep{hou2024llmse,zhang2024unifying,jin2024llmagents,wang2024llmagents,liu2024llmagents},
a process that produced 17 documented borderline cases and a decision
flow chart with eight sequential questions.
Each screening decision was recorded in a versioned spreadsheet with
per-criterion rationale, permitting retrospective audit of the
exclusions.

Data extraction was performed by the same researcher, introducing risk
of interpretation bias in qualitative fields such as architecture-type
classification and degradation mechanism.
Four mitigations reduced this risk: (i) programmatic validation of 448
enum values (16 fields $\times$ 28 studies), with error rate below 1\%
after corrections; (ii) spot-check of 40 interpretive fields across 10
selected studies, as detailed in \autoref{sec:extracao}; (iii) factual
fidelity digests generated from the full text of each study, used as
reference when drafting the manuscript's interpretive claims; and (iv)
re-verification of the \textit{negative\_result} fields against the
original PDF.
Even with these mitigations, formal reproducibility of the
classification in interpretive fields was not assessed.

The preprint prevalence (82\% of the corpus) reduces confidence in this
review's findings.
The field of knowledge persistence for coding agents is recent, with
publications concentrated between 2024 and 2026, and only five studies
have been published in peer-reviewed venues: two at
conferences~\citep{wang2026improving,zhang2025darwingodel} and three at
workshops~\citep{deshpande2025memtrack,shen2025talm,lindenbauer2025ctimrover}.
Sensitivity analysis S1 (\autoref{sec:sensibilidade}) shows that,
when preprints are excluded, no qualitative conclusion inverts, but the
evidence base for cost (SQ4) drops from seven to two studies and the
detailed taxonomy loses empirical support.

Searches were limited to two databases (Scopus and arXiv), complemented
by citation chasing.
Studies published exclusively in uncovered bases (for example, IEEE
Xplore, ACM Digital Library) may have been missed.
The rationale for this choice, based on \citet{mourao2020hybrid}, is
documented in \autoref{sec:fontes}.

The 14 studies that use SWE-bench Verified differ in base model, agent
framework, instance subset, and run conditions (\autoref{sec:qs2}).
The construct ``agent performance'' is not operationalised uniformly
across studies, which limits direct comparability of resolution
percentages reported as if they were equivalent quantities.
The taxonomy proposed in this review for persistence architectures also
embeds interpretive choices: families such as ``hybrid'' group
heterogeneous mechanisms whose fine differentiation would require
metadata absent from published texts.

Heterogeneity across studies precludes formal meta-analysis and limits
inferences to qualitative patterns synthesised narratively.
Classifying findings into evidence types A, B, and C
(\autoref{tab:certeza}) makes the empirical support of each inference
explicit: type A for direct counts, type B for aggregated cross-study
patterns, and type C for mechanism characterisations.
The two sensitivity analyses conducted (S1, preprint exclusion;
S2, exclusion of studies scoring below 3.0) did not invert any
quantitative finding, reinforcing the stability of the synthesis under
alternative inclusion criteria.

\subsection{Research agenda}
\label{sec:agenda-pesquisa}

The three gaps confirmed by this review define the field's immediate
agenda.
Gap~1 --- the absence of head-to-head comparisons --- calls for an
experiment that evaluates at least three representative architectures
on the same benchmark, with the same model, and under the same compute
budget.
Gap~2 --- cost underreporting --- requires that such an experiment
include an integrated cost-per-task analysis, producing an explicit
Pareto frontier between accuracy and cost.

Metric standardisation is a prerequisite for future comparability.
The continual-learning metrics proposed by
\citet{joshi2025swebenchcl} (\textit{CL-Score},
\textit{forward/backward transfer}, \textit{forgetting}) have been
adopted by no other study in the corpus.
Adopting standardised cost metrics (tokens per task, monetary cost,
wall-clock time) and knowledge-retention metrics would enable
cross-study comparisons that are currently unfeasible.

Only one of the 24 primary studies reports memory sharing across
agents~\citep{chen2025sweexp}, and even there it is limited to
experience reuse within the same framework.
Sharing protocols across heterogeneous agents remain without empirical
evaluation.
The vulnerabilities flagged by \citet{srivastava2025memorygraft}
indicate that any such protocol must address the risk of poisoning in
open stores before it can be recommended for production.
